% Appendix: implementation. Values are those of the configuration used for the reported build.
\section{Implementation Details}
\label{app:implementation}

\begin{algorithm*}[t]
\caption{Assembling the page of entity $e$ at cutoff $t$.}
\label{alg:assemble}
\small
\begin{algorithmic}[1]
\Require graph of verifier-accepted intervals; entity $e$; cutoff $t$
\State $I \gets$ copies of $\{\, i \in \textsc{Intervals}(e) : i.\mathit{start} \le t \,\}$ \Comment{a local view; the graph is not modified}
\If{$I = \varnothing$} \Return no page \Comment{an entity with no interval by $t$ has no page and is no link target}
\EndIf
\For{$i \in I$}
  \State $i.\mathit{quotes} \gets \{\, q \in i.\mathit{quotes} : \mathrm{vol}(q) \le t \,\}$ \Comment{each quote filtered by its own volume}
  \If{$i.\mathit{end} > t$} \State $i.\mathit{end} \gets \bot$ \Comment{do not reveal that, or when, it ends}
  \EndIf
\EndFor
\State $I \gets \{\, i \in I : \textsc{Verified}(i) \,\}$ \Comment{facts the verifier rejected are withheld}
\State $\mathit{fields} \gets \textsc{Render}(I)$ \Comment{deterministic: names, attributes, relations, history}
\State $\mathit{prose} \gets \textsc{Generate}(I)$ \Comment{model sees only $I$; no evidence $\Rightarrow$ no call, empty section}
\State $\mathit{page} \gets \textsc{Link}(\mathit{fields}, \mathit{prose}, \{\, s : \mathrm{vol}(s) \le t \,\})$ \Comment{surface forms and targets revealed by $t$}
\State \Return $\mathit{page}$ with cutoff-$t$ title and navigation
\end{algorithmic}
\end{algorithm*}

\paragraph{Claim schema.} A claim is a record with a subject (an entity identifier, never a
raw surface form), a predicate, a kind (attribute, relation or trait), an object entity or a
scalar value, an optional short qualifier, a polarity (asserted, denied or presumed), the
volume and chapter of its first evidence, and one or more evidence items, each a paragraph
identifier with a verbatim quote of at most 240 characters. \emph{Presumed} is meant to keep a belief
apart from a fact: a character thought dead in one volume is recorded as presumed dead there,
not as dead. The renderer marks a presumed value in attribute and relation lines, but the
infobox's family field and the list of related characters print it unmarked
(Section~\ref{sec:results-main}). Polarity and qualifier are part of a claim's identity, so two observations that
disagree on either never merge. A claim whose subject or object cannot be resolved to a known
entity is dropped.

\paragraph{Ontology.} Entities are typed as character, group, place, event or object; only
characters get their own pages, and the other types share codex pages. There are nine
attributes (age, gender, status, rank, title, nickname, origin, occupation, appearance), five
trait keys (personality, motivation, skill, fear, background) and twenty relation predicates
covering kinship (parent, child, sibling, relative, spouse), affect (romantic, friend, rival,
enemy), roles (mentor, student, commands, serves under, comrade, affiliated with) and events
(origin, participated in, pilots, killed, saved). The ontology is shared by all series and was
not adapted to the evaluated work.

\paragraph{Extraction settings.} The chapter pass reads spans of about 4{,}000 tokens with 500
tokens of overlap and records, for each scene, a summary, the characters who take part and up
to four paragraphs that support the summary. The mention pass reads a window of four
paragraphs before and after a mention, never across a hard scene break, with at most 40
windows per entity per volume, one window per call and at most 6{,}000 tokens per window (a
denser cluster of mentions is split into consecutive windows). The long-context pass reads all of a character's paragraphs in a volume
(up to 150{,}000 tokens) for gender, age and origin, and separately for background. A claim
needs a self-reported confidence of at least 0.5 to be kept and 0.7 to fill a structured
field. An evidence-adequacy gate rejects quotes shorter than six words, quotes in which the
subject is not named within one paragraph, and quotes carrying irrealis markers (``would
have'', ``might have'').

\paragraph{Intervals.} Single-valued attributes supersede: a value's interval closes when the
next distinct value first appears. Relations in a declared conflict group supersede the same
way; every other relation and multi-valued attribute is an open interval that opens at its
earliest evidence, with its qualifier taken from evidence no later than that volume. Free-text
values whose BGE-M3 embeddings have cosine similarity of at least 0.86 are clustered before
intervals are built.

\paragraph{Generation.} Every model call uses \res{GEN-model}. The intended character-page
template specifies appearance 2--4, personality 2--6, and history 3--12 sentences; shared-page
summaries request one sentence. The archived \sys appearance prompt instead requests
3--4 sentences (Appendix~\ref{app:prompts}). These are prompt requirements, not a measurement
of output-length enforcement. Generation prompts request keys identifying the input lines
supporting the text, which map back to source paragraphs. The parser also normalizes returned
paragraph IDs during citation resolution, as in the codex example below. A shared-page
summary is retained only if a second call judges every statement to be established by its
input passages. A rejected summary is regenerated once and omitted if it fails again.
Responses are cached by provider, model, prompt, and parameters, so rebuilding an unchanged
configuration requires no new calls. All prompts are released verbatim with the code.
Prompts for \sys and text-reading comparison systems do not state the work's title or
author, but mined names can carry it: the series title is an entity in \sys's extraction
prompts, and the typing model's exhibited output names the novel (Appendix~\ref{app:prompts-full}).
The closed-book diagnostic receives the title by design.
A character with no STATUS claim is printed as \emph{alive (assumed)}: a configured
default with no evidence behind it, marked as assumed on the page, since the absence of a
death claim does not establish that a character is alive.

\paragraph{Prompt and response defects.} The extraction prompt of the reported runs
exposes internal revision labels, module paths, and implementation notes in the STATUS and
APPEARANCE descriptions. The prompt builder passes the same canonicalization note to the
verifier as the \emph{definition} of APPEARANCE, which the verifier is told to judge
against; a rule for deduplicating values does not define a physical characteristic. The
exhibited call accepts its appearance claim; the effect on other decisions is not measured.
The codex response returns a paragraph ID where the prompt requests an input-line key; the
parser maps it to its paragraph during citation resolution, without checking that the
paragraph supports the sentence. All reported results use this configuration; a repaired
prompt builder would require new runs, which this paper does not report.

\paragraph{Changes between builds.} The reported build differs from the first
build (Section~\ref{sec:results-quality}, Table~\ref{tab:builds}) by one batch of changes,
made after the first build had been labeled and applied to every series alike.
\emph{Entity resolution:} a bare name and a fuller name first seen in a later volume go to
the model instead of merging for free; a first name shared by several people is assigned per
volume; fuzzy matching may not exchange a whole word; the cap on model-judged name pairs
rose from 1{,}000 to 5{,}000; hyphenated names and titles of the form ``X of the Y'' are
mined whole. \emph{Ingestion:} publisher's back matter is dropped. \emph{Scenes:} the
participants of a scene are those who take part, not those mentioned, and each summary names
up to four supporting paragraphs. \emph{Citations:} generated text names the input lines it
rests on, and scene lines print their paragraphs on timeline, pair and codex pages; the first
build cited a chapter for a scene line and nothing for most prose. \emph{Gates:} a prose
answer that says it found nothing is discarded, and summaries on shared pages pass the check
described above. \emph{Windows:} the 6{,}000-token ceiling on a mention window, which the
first build did not enforce on a dense cluster. One side effect: the first
build held two entities for several main characters, and each entity received its own
windows, so the same text was read twice; with the entities merged it is read once, and the
claims stored for two such characters fell from 162 and 66 to 118 and 46 on as many windows.
We did not isolate this change's contribution to the reported build's lower recall.
